\documentclass[../../main.tex]{subfiles}

\begin{document}

\section{极大原理}\label{section:ch01-极大原理}

我们已经说过，由\cref{axiom:ch01-选择公理}可以得出任何集合都能良序化这一深刻的定理。在数学中，\cref{axiom:ch01-选择公理}还有更加重要的推论。这里所提及的“极大原理”有多种版本。在 1914 年至 1935 年间，多位数学家曾对极大原理独立地予以论述，他们包括 F. Hausdorff、K. Kuratowski、S. Bochner 及 M. Zorn。而当时一种典型的处理思路是用良序原理证明他们各自的“极大原理”。后来人们发现，实际上这些“极大原理”都是与良序原理等价的。以下我们就其中的一些予以介绍。

首先给出一个定义。给定集合 $A$，$A$ 中的一个关系 $\prec$ 称为 $A$ 上的严格偏序（strict partial order），如果它具有下面两个性质：
\begin{enumerate}[(1)]
\item （非自反性）关系 $a \prec a$ 不成立.
\item （传递性）若 $a \prec b$ 并且 $b \prec c$，则 $a \prec c$.
\end{enumerate}
这刚好是全序性质（见第\cref{section:ch01-关系}节）的第二条和第三条，去掉了可比较性。换言之，一个严格偏序关系与全序关系十分类似，只是不要求对于集合中两个不同的点 $x$ 和 $y$，或者 $x \prec y$ 成立，或者 $y \prec x$ 成立。

如果 $\prec$ 是集合 $A$ 上的一个严格偏序关系，则不难看出 $A$ 有某子集 $B$，它是这个关系下的全序集，也就是说 $B$ 中任意元素偶对在 $\prec$ 下都可以比较。

现在讲述以下极大原理，它是 1914 年由 Hausdorff 首先提出的。

\begin{theorem}[极大原理]\label{theorem:ch01-极大原理}
设 $A$ 是一个集合，$\prec$ 为 $A$ 上一个严格偏序。则存在 $A$ 中的一个极大全序子集 $B$。

换句话说，存在 $A$ 中的子集 $B$，使得 $B$ 是一个关于 $\prec$ 的全序集，并且 $A$ 中任何以 $B$ 作为其真子集的集合都不是关于 $\prec$ 的全序集。
\end{theorem}

\par
\begin{figure}[htbp]
\centering
\includegraphics[width=0.6\textwidth]{ch01-fig11-1}
\caption{平面圆域族的两个极大全序子族}\label{figure:ch01-极大圆域}
\end{figure}

\begin{example}\label{example:ch01-极大圆域}
如果 $\mathcal{A}$ 是任意一个集族，关系“是某集合的一个真子集”为 $\mathcal{A}$ 上的一个严格偏序关系。设 $\mathcal{A}$ 为平面上所有圆域（圆周的内部）的族。$\mathcal{A}$ 的一个极大全序子族由所有中心在原点的圆域组成。另一个极大全序子族由所有下述圆域组成，其边界圆周在 $y$ 轴右边并且在原点处与 $y$ 轴相切，参见\cref{figure:ch01-极大圆域}。

\end{example}

\begin{example}\label{example:ch01-水平直线}
若 $(x_{0}, y_{0})$ 和 $(x_{1}, y_{1})$ 为平面 $\mathbb{R}^{2}$ 中的两点，当 $y_{0} = y_{1}$ 并且 $x_{0} < x_{1}$ 时，定义
\begin{align*}
(x_{0}, y_{0}) \prec (x_{1}, y_{1}).
\end{align*}
这是 $\mathbb{R}^{2}$ 上的一个偏序，在这个关系下，两个点仅当它们位于同一水平线上时才可以比较。于是极大全序集为 $\mathbb{R}^{2}$ 中的水平直线。
\end{example}

我们只能给出极大原理的一个直观“证明”。把它分为几步来做，下面予以形象地描述。取一个盒子，我们依以下方式把 $A$ 的某些元素依次放入盒中：首先任意选取 $A$ 的一个元素，将其放入盒子中。然后，在 $A$ 中另选一个元素，如果它与盒子中的元素都可以比较，则将其放入盒子中，否则，就将其放在一边。到某一步，我们就有了在盒中元素的一个族以及被扔在一边的元素的族。再在 $A$ 剩余元素中取一个，如果它与盒中每一个元素都可以比较，便把它放入盒子中，否则，将其放在一边。以此方式一直做下去。当你对 $A$ 中的所有元素都验证之后，则盒子中的元素是可以互相比较的，因而它们组成一个全序集。而不在盒中的每一个元素都至少与盒中的一个元素不能比较，这正是我们把它放在一边的原因。由于 $A$ 已经不再有满足可比条件的更大的子集，因此，盒中的全序集是极大的。

当然，上述“证明”有一个薄弱环节，就是我们说的“当你对于 $A$ 的所有元素都验证之后”。你怎么知道你“做完了”对于 $A$ 中所有元素的验证呢？如果 $A$ 是可数集，那不难把这种直观证明变成实际的证明。我们来看可数无限的情形（有限情形是比较容易的）。将 $A$ 的元素逐一给以各不相同的正整数指标，即
\begin{align*}
A = \{a_{1}, a_{2}, \cdots\}.
\end{align*}
这种指标提供了如何验证 $A$ 中元素的方法，也可以使我们知道什么时候全部验证完毕了。

让我们定义一个函数 $h: \mathbb{Z}_{+} \longrightarrow \{0, 1\}$，它的指派法则是：“将 $a_{i}$ 放入盒中”的 $i$ 对应于 $0$，“将 $a_{i}$ 放在一边”的 $i$ 对应于 $1$。这意味着：$h(1) = 0$。对于每一个 $i > 1$，我们有 $h(i) = 0$ 当且仅当 $a_{i}$ 与集合
\begin{align*}
\{a_{j} \mid j < i \text{ 和 } h(j) = 0\}
\end{align*}
中的每一元素都是可比较的。由\cref{theorem:ch01-归纳定义原理}，这个公式决定了唯一函数 $h: \mathbb{Z}_{+} \longrightarrow \{0, 1\}$。容易验证：所有使得 $h(i) = 0$ 的 $a_{i}$ 的集合是 $A$ 的一个极大全序子集。

如果 $A$ 不是可数集，只要允许我们使用良序定理，就可以将上述程序改变得可行。我们以某一良序集 $J$ 代替原来的 $\mathbb{Z}_{+}$，对 $A$ 的元素逐一加以标记（即建立一一对应），使得 $A = \{a_{\alpha} \mid \alpha \in J\}$。为此，我们需要良序定理，得到 $A$ 与良序集 $J$ 之间有一个一一对应。这时便可以像上一段中那样，用 $\alpha$ 代替 $i$ 进行论证。严格地讲，需要把归纳定义原理推广到良序集上去，而这件事并不是十分困难（见附加习题）。

因此，良序定理蕴涵着极大原理。

尽管 Hausdorff 极大原理出现得最早，并且可能是最易于理解的一种说法，然而，当今引用最多的却是另一个与之有关的说法。这就是人们常说的 Zorn 引理，尽管 Kuratowski（1922）及 Bochner（1922）都曾先于 Zorn（1935）阐述并证明了某种形式的 Zorn 引理。相关历史以及与之有关的资料，请参见文献 $[\mathrm{C}]$ 及 $[\mathrm{Mo}]$。为陈述这一原理，我们需要引进某些术语。

\begin{definition}[上界与极大元]\label{definition:ch01-上界极大元}
设 $A$ 是一个集合，$\prec$ 是 $A$ 上的一个严格偏序。若 $B$ 为 $A$ 的一个子集，$c \in A$，使得对于每一个 $b \in B$ 都有 $b \prec c$ 或者 $b = c$，则称 $c$ 是 $B$ 的一个上界（upper bound）。若 $m$ 是 $A$ 的一个元素，并且对于每一个 $a \in A$，$m \prec a$ 都不成立，则称 $m$ 是 $A$ 的一个极大元（maximal element）。
\end{definition}

\begin{axiom}[Zorn 引理]\label{axiom:ch01-Zorn}
设 $A$ 是一个严格偏序集。若 $A$ 的每一个全序子集有上界，则 $A$ 中必有一个极大元。
\end{axiom}

Zorn 引理是极大原理的一个简单推论：给定 $A$，极大原理蕴涵着 $A$ 有极大的全序子集 $B$。由 Zorn 引理的假设可见，$B$ 在 $A$ 中有上界 $c$。$c$ 自然是 $A$ 的一个极大元。事实上，若存在某 $d \in A$ 使得 $c \prec d$，那么对于每一个 $b \in B$ 有 $b \prec d$，从而，$B \cup \{d\}$ 是一个全序集且以 $B$ 为其真子集，这与 $B$ 极大性矛盾。

事实上，极大原理也是 Zorn 引理的一个简单推论。见习题 5$\sim$7。

最后还要加一个附注。我们已经给出了集合上严格偏序的定义，却没有说偏序是什么。设 $\prec$ 是集合 $A$ 上的一个严格偏序。将 $a \leqslant b$ 定义为 $a \prec b$ 或者 $a = b$，则关系 $\leqslant$ 为 $A$ 上的一个偏序（partial order）。例如，集族中的包含关系 $\subset$ 是一个偏序，而真包含关系是一个严格偏序。

大多数作者宁可讨论偏序，而不讨论严格偏序。极大原理和 Zorn 引理也往往用偏序加以描述。究竟使用哪一种，这仅是一个嗜好与方便的问题。

\subsection*{习题}

\begin{exercise}\label{exercise:ch01-极大原理-1}
若 $a$ 和 $b$ 是两个实数，当 $b - a$ 为正有理数时定义 $a \prec b$。证明这是 $\mathbb{R}$ 上的一个严格偏序。它的极大全序子集是什么？
\end{exercise}

\begin{exercise}\label{exercise:ch01-极大原理-2}
(a) 设 $\prec$ 是集合 $A$ 中的一个严格偏序。$A$ 中的一个关系 $\leqslant$ 定义为：当 $a \prec b$ 或者 $a = b$ 时，$a \leqslant b$。证明这个关系具有下述称之为偏序公理（partial order axioms）的性质：
\begin{enumerate}[(i)]
\item 对于所有 $a \in A$，$a \leqslant a$.
\item $a \leqslant b$ 和 $b \leqslant a \Longrightarrow a = b$.
\item $a \leqslant b$ 和 $b \leqslant c \Longrightarrow a \leqslant c$.
\end{enumerate}

(b) 设 $P$ 为 $A$ 上满足性质 (i)、(ii)、(iii) 的一个关系。$A$ 上的关系 $S$ 定义为：当 $a P b$ 并且 $a \neq b$ 时，$a S b$。证明 $S$ 是 $A$ 上的一个严格偏序。
\end{exercise}

\begin{exercise}\label{exercise:ch01-极大原理-3}
设 $A$ 是有严格偏序 $\prec$ 的一个集合，$x \in A$。如果我们想找出 $A$ 的包含 $x$ 的一个极大全序子集 $B$，一个办法是试图将 $B$ 定义为 $A$ 中所有可与 $x$ 比较的那些元素的集合，即
\begin{align*}
B = \{y \mid y \in A \text{ 并且或者 } x \prec y \text{ 或者 } y \prec x\}.
\end{align*}
但是这并不总是可行的。\cref{example:ch01-极大圆域}和\cref{example:ch01-水平直线}中的哪一个是按照这个方式去做的，哪一个不是？
\end{exercise}

\begin{exercise}\label{exercise:ch01-极大原理-4}
给定 $\mathbb{R}^{2}$ 中的两点 $(x_{0}, y_{0})$ 和 $(x_{1}, y_{1})$。当 $x_{0} < x_{1}$ 并且 $y_{0} \leqslant y_{1}$ 时定义
\begin{align*}
(x_{0}, y_{0}) \prec (x_{1}, y_{1}).
\end{align*}
证明：曲线 $y = x^{3}$ 与 $y = 2$ 都是 $\mathbb{R}^{2}$ 的极大全序子集。而曲线 $y = x^{2}$ 不是。试求出所有极大全序子集。
\end{exercise}

\begin{exercise}\label{exercise:ch01-极大原理-5}
证明：Zorn 引理蕴涵着以下引理：

\begin{lemma}[Kuratowski 引理]\label{lemma:ch01-Kuratowski}
设 $\mathcal{A}$ 为一个集族。如果 $\mathcal{A}$ 的每一个对于真包含关系而言的全序子族 $\mathcal{B}$ 的所有元素之并仍属于 $\mathcal{A}$，则 $\mathcal{A}$ 中有一个元素不真包含于 $\mathcal{A}$ 的任何元素之中。
\end{lemma}
\end{exercise}

\begin{exercise}\label{exercise:ch01-极大原理-6}
设 $\mathcal{A}$ 为集合 $X$ 的子集的一个族。若 $X$ 的子集 $B$ 属于 $\mathcal{A}$ 当且仅当 $B$ 的每一个有限子集属于 $\mathcal{A}$，则称 $\mathcal{A}$ 是具有有限特征的。证明 Kuratowski 引理蕴涵着以下引理：

\begin{lemma}[Tukey 引理]\label{lemma:ch01-Tukey}
设 $\mathcal{A}$ 是一个集族。若 $\mathcal{A}$ 是具有有限特征的，则 $\mathcal{A}$ 中存在一个元素，使得它不真包含于 $\mathcal{A}$ 的任何元素之中。
\end{lemma}
\end{exercise}

\begin{exercise}\label{exercise:ch01-极大原理-7}
证明 Tukey 引理蕴涵着 Hausdorff 极大原理。[提示：设 $\prec$ 是 $A$ 上的一个严格偏序。用 $\mathcal{A}$ 表示按照相对于 $\prec$ 而言 $A$ 的所有全序子集的族。证明 $\mathcal{A}$ 是具有有限特征的。]
\end{exercise}

\begin{exercise}\label{exercise:ch01-极大原理-8}
Zorn 引理在代数中的一个典型应用是证明每一个向量空间有一个基。首先回忆几个有关概念。设 $A$ 为向量空间的一个子集，倘若一个向量是 $A$ 的某些元素的有限线性组合，则称它是一个由 $A$ 所张成的向量。集合 $A$ 被称为是无关的，若将零向量写成 $A$ 中有限个元素的线性组合时，所有系数必全为零。如果 $A$ 是无关的并且 $V$ 的每一向量都是由 $A$ 所张成的向量，那么 $A$ 是 $V$ 的一个基。
\begin{enumerate}[(a)]
\item 若 $A$ 是无关的并且 $v \in V$ 不是由 $A$ 所张成的，证明 $A \cup \{v\}$ 是无关的.
\item 证明 $V$ 的所有无关的子集的族有极大元.
\item 证明 $V$ 有一个基.
\end{enumerate}
\end{exercise}

\end{document}